\documentclass[11pt,letterpaper]{article}
\usepackage[T1]{fontenc}
\usepackage{lmodern}
\usepackage[margin=1in]{geometry}
\usepackage{amsmath,amssymb,amsthm,mathtools}
\usepackage{booktabs,longtable,array,enumitem}
\usepackage{microtype,xcolor,xurl}
\usepackage[colorlinks=true,linkcolor=blue!45!black,urlcolor=blue!45!black,citecolor=blue!45!black]{hyperref}
\usepackage{fancyhdr}
\pagestyle{fancy}
\fancyhf{}
\fancyhead[L]{\small Residue orientation and root transfer}
\fancyhead[R]{\small ProveIt research contribution}
\fancyfoot[C]{\thepage}
\setlength{\headheight}{14pt}
\setlength{\emergencystretch}{2em}
\setlength{\parskip}{4pt}
\setcounter{tocdepth}{2}
\newtheorem{theorem}{Theorem}[section]
\newtheorem{lemma}[theorem]{Lemma}
\newtheorem{proposition}[theorem]{Proposition}
\newtheorem{corollary}[theorem]{Corollary}
\theoremstyle{remark}
\newtheorem{remark}[theorem]{Remark}
\newcommand{\Z}{\mathbb Z}
\newcommand{\N}{\mathbb N}
\newcommand{\Con}{\operatorname{Con}}
\newcommand{\sgn}{\operatorname{sgn}}
\newcommand{\pos}{\operatorname{Pos}}
\hypersetup{pdftitle={Residue Orientation and Root Transfer for Diophantine Consistency Equations},pdfauthor={AI-assisted research prepared for Vladimir Reshetnikov},pdfsubject={Explicit polynomial equations equivalent to syntactic inconsistency of ZFC}}
\title{\textbf{Residue Orientation and Root Transfer}\\[5pt]
\large Smaller explicit Diophantine equations for syntactic consistency of ZFC}
\author{AI-assisted research prepared for Vladimir Reshetnikov\\
\small A contribution for the ProveIt repository}
\date{23 October 2026}
\begin{document}
\maketitle
\thispagestyle{empty}
\begin{abstract}
We give two proof-driven reductions of explicit Diophantine consistency
equations.  First, a signed Pell argument recovers a coordinate as
$\pm2^R$ with $R$ odd; affine congruences then select the positive sign.
This removes a dedicated square witness and gives a quartic in forty
unrestricted integer variables.  Exact coordinate eliminations give a
sextic in thirty-four and an octic in twenty-nine integer variables.
Second, we transfer Pell roots into a transport gap.  A two-root norm
gives twelve strictly positive variables of degree eighty-four in 103
arithmetic operations.  A three-root extension uses a coprimality argument
to recover an auxiliary integer root, giving eleven positive variables
of degree 325 in 113 operations.  A thirteen-positive-variable companion
has degree thirty in 94 operations.
These improve specified points in the inspected repository snapshot.
Every fixed coefficient is supplied exactly, with an offline verification
recipe.  The new equivalences are proved by ordinary mathematics using
explicitly named inherited compiler, Pell, and syntactic calibration
theorems.  They are not new end-to-end proof-assistant certificates, and
no global minimality or proof of consistency is asserted.
\end{abstract}
\vfill
\noindent\textbf{Source snapshot.}\quad
\texttt{9cca523d9f27ad39d52993f4065cbd7e58f8bdd4}\\
\noindent\textbf{Primary deliverables.}\\
\texttt{integer/quartic40\_direct.json}\\
\texttt{eleven/p11\_positive\_degree325.json}\\
\texttt{positive/p12\_positive\_degree84.json}\\
\noindent\textbf{Reproduction.}\quad
\texttt{python3 verify.py}; \texttt{python3 verify.py --optimized}.
\clearpage
\tableofcontents
\clearpage

\section{Results, comparison, and the meaning of simplicity}
\label{sec:results}

Fix the repository's recursive syntactic proof predicate for ZFC.  In this
article $\Con(\mathrm{ZFC})$ means that no finite derivation of falsity
exists for that predicate.  The target is an explicitly specified
polynomial $P\in\Z[x_1,\ldots,x_m]$, together with a stated domain $D$, for
which ordinary mathematics formalizable in ZFC proves
\begin{equation}
 \bigl(\exists \mathbf x\in D^m\;P(\mathbf x)=0\bigr)
 \quad\Longleftrightarrow\quad
 \text{there is a finite syntactic ZFC refutation}.
 \label{eq:target}
\end{equation}
Negating the two sides yields exactly the requested equivalence between
unsolvability and $\Con(\mathrm{ZFC})$.  This establishes a relationship
between two statements; it does not establish their truth or falsity.

The phrase ``simplest equation'' has several competing meanings.  A
reduction in variables can increase degree, and treating a large fixed
integer as one arithmetic-circuit leaf can hide most of the written
description.  We therefore report witness domain and count, exact total
degree, number and degree of residuals, arithmetic operations when
optimized, and fixed coefficient storage separately.  A named intermediate
register is a polynomial abbreviation, never an additional quantified
unknown.

\begin{table}[ht]
\centering
\small
\begin{tabular}{@{}llrrl@{}}
\toprule
Equation & Witness domain & Variables & Degree & Binary operations\\
\midrule
$Q_{40}^{\mathrm{dir}}$ & $\Z$ & 40 & 4 & Not optimized\\
$Q_{34}$ & $\Z$ & 34 & 6 & Not optimized\\
$Q_{29}$ & $\Z$ & 29 & 8 & Not optimized\\
$P_{11}$ & $\Z_{>0}$ & 11 & 325 & 113\\
$P_{12}$ & $\Z_{>0}$ & 12 & 84 & 103\\
$P_{13}$ & $\Z_{>0}$ & 13 & 30 & 94\\
\bottomrule
\end{tabular}
\caption{New single polynomial equations.  The integer equations are
sums of 25, 19, and 14 squared residuals, respectively.}
\label{tab:new}
\end{table}

\begin{theorem}[The delivered consistency equations]
Each polynomial in Table~\ref{tab:new}, with its own complete fixed
coefficient tuple and stated witness domain, satisfies
\eqref{eq:target}.  Its exact degree and number of variables are as
listed.  The claims use the inherited source theorems identified in the
proofs and provenance files; they do not introduce a consistency
assumption.
\end{theorem}

The integer constructions improve the inspected degree-four, degree-six,
and degree-eight counts from $41,35,30$ to $40,34,29$.  The parent quartic
is package 64's $Q_{41}$; the two smaller-count comparison companions are
package 62's sextic and octic.  The twelve-positive construction improves
package 61's degree-$118$, 104-operation equation in both measures:
$118\to84$ is a $28.81\%$ degree reduction, and $104\to103$ saves one
operation.  The thirteen-positive comparison is a tradeoff:
degree $34\to30$, with operations $93\to94$.  These are comparisons to
literal artifacts in the pinned repository, not just to the collection
README, which still described an earlier wave at the inspection time.
Relevant comparison files and their hashes are retained in
\texttt{baseline/} and \texttt{BASELINE.json}.

The eleven-positive extension improves package 52's degree-$954$,
113-operation point to degree $325$ at the same 113-operation cost.
This is a $65.93\%$ degree reduction.  Its two residuals have degrees
$312$ and $13$.  A modulo-four unit-product argument combines them
without taking their sum of squares; the final polynomial has degree
$312+13=325$.

\begin{table}[ht]
\centering
\small
\begin{tabular}{@{}lll>{\raggedright\arraybackslash}p{4.0cm}@{}}
\toprule
Slice & Inspected predecessor & New result & Nature of the change\\
\midrule
Quartic over $\Z$ & 41 variables, package 64 & 40 & One sign witness removed\\
Sextic over $\Z$ & 35 variables, package 62 & 34 & Same orientation, exact elimination\\
Octic over $\Z$ & 30 variables, package 62 & 29 & Same orientation, exact elimination\\
Eleven positive & Degree 954 / 113 ops, package 52 & 325 / 113 & Three roots in one transport carrier\\
Twelve positive & Degree 118 / 104 ops, package 61 & 84 / 103 & A strong root is eliminated\\
Thirteen positive & Degree 34 / 93 ops, package 61 & 30 / 94 & Adaptive modulus composition\\
\bottomrule
\end{tabular}
\caption{A bounded comparison, with domains and costs kept distinct.}
\end{table}

The small-variable high-degree records remain separate: the repository
already contains nine-positive and ten-unrestricted-integer constructions.
The present work does not lower those counts.  Likewise, the fifteen-positive
degree-$24$ and 66-operation constructions address different objectives.
The recommended $Q_{40}^{\mathrm{dir}}$ has 393 expanded symbolic monomials, compared
with 388 in its parent; the alternate projection parent has 614.
No coefficient-height, expanded-length, or global circuit-cost reduction
is claimed for that family.  A mathematically new mechanism can improve
one boundary of the comparison without dominating every other equation.

\subsection{What is new in this iteration}

The first method changes the order in which signs are established.  A
positive-square guard is usually imposed before a Pell argument.  Here
the argument is proved with signed $X,Y,k,c$ throughout, and gives
$X=\pm2^R$ with $R$ odd.  The already imposed congruence
$X\equiv2\pmod3$ then rules out the negative branch.  The congruence is
enforced most directly by $X_{\rm old}=3X_3-1$, so it needs no
new witness and preserves quadratic residual degree.  A companion
version restricts the two existing scale coordinates instead, retaining
a monic $X$ definition for the sextic and octic projections.

The second method exploits an existing free positive transport gap.
Instead of keeping a strong root $f$ and a gap $r$ separately, use the
coordinate $n=r-L+f$, with $L>|f|$, and eliminate $f$ by a norm.  The
proof must handle both signed roots recovered by that norm.  A stronger
transport multiplier guarantees the inverse map has $n>0$.  The
remaining auxiliary relations use only $f^2$, so their degree remains
controlled after substitution.

The eleven-variable extension uses that normalization to make the three
radicands independent of the signs of their roots.  It retains the
adaptive main Pell equation separately and puts the first, strong, and
auxiliary roots into the transport carrier.  Multiplying the normalized
strong ordinate by $q(q-1)$ makes the auxiliary denominator coprime to
the transport denominator, so a rational root recovered from the norm
must be an integer.  Elementary separation bounds exclude every
irrational cancellation.  Norms, positive majorants, and unit-product
finalizers already appear in earlier packages; the new claim is this
particular arithmetic composition and its measured endpoint.

Both mechanisms are new relative to the inspected source packages.
We make no literature-wide priority claim.  The thirteen-variable
adaptive carrier and the low-degree coordinate eliminations explicitly
reuse earlier ideas; their attribution is retained in the component
notices.

\subsection{Why the final constructions retain the native endpoint}

The user permitted a fresh computational or logical source.  We examined
direct provability encodings as well.  Carl and Moroz give an explicit
Diophantine treatment of provability in predicate calculus and discuss
set-theoretic applications~\cite{carlmoroz}.  Formisano, Gavazzi and
Omodeo investigate a relational equational presentation of NBG; their
2025 article describes preliminary work on its proper-axiom
component~\cite{fgo}.  These are useful routes to smaller transparent
proof checkers, but our inspection did not produce a fully materialized
alternative with a competitive, proved coefficient and variable count.
An uncompiled proof checker would not be a replacement for the requested
explicit equation.

The delivered constructions therefore retain the existing native
K215-to-polynomial source and improve its arithmetic representation.
They do not use Jones's universal polynomial or Sun's ten-integer
finalizer.  Jones's work~\cite{jones} and Sun's recent ten-variable
theorem~\cite{sun} explain relevant comparison boundaries, but are not
newly invoked to justify these polynomial equations.  Native
compiler semantics and K215 calibration remain explicit inherited
dependencies.

\section{Residue orientation: a quartic in forty unrestricted integers}
\label{int:main}

The main structural improvement of this contribution is the removal of a
dedicated positivity witness.  Source~64 uses the equation $X=2u^2$ to
start its Pell argument with $X\geq0$.  We allow $X$ to have either sign,
recover $X=\pm2^R$ with $R$ odd, and select the positive branch by the
congruence $X\equiv2\pmod3$.  The congruence is obtained by affine
restrictions on two existing coordinates, so it introduces no unknown.
The preferred quartic uses the even simpler direct restriction
$X=3X_3-1$ and has $393$ expanded symbolic monomials.  We first prove the
equivalent $40$-variable version with two affine restrictions, whose monic
$X$ definition supports the best sextic and octic projections; the direct
version is specified in Subsection~\ref{int:direct}.

\begin{theorem}\label{int:theorem}
There is an explicitly specified fixed polynomial $Q_{40}^{\mathrm{aff}}$ in forty
unrestricted integer unknowns, of exact total degree four, such that
ordinary mathematics formalizable in ZFC proves
\[
  (\exists\mathbf x\in\mathbb Z^{40})\ Q_{40}^{\mathrm{aff}}(\mathbf x)=0
  \quad\Longleftrightarrow\quad
  \text{there is a finite syntactic ZFC refutation}.
\]
Consequently ZFC proves that the absence of integer zeros of $Q_{40}^{\mathrm{aff}}$
is equivalent to $\operatorname{Con}(\mathrm{ZFC})$.
Exact polynomial projections give corresponding equations in thirty-four
integers of degree six and twenty-nine integers of degree eight.
\end{theorem}

The statement imports the explicitly identified native-compiler, fixed-tuple,
K215 calibration, and Pell-rank theorems of the repository.  The new part is
the sign-uniform power-recovery argument below and its application to literal
polynomials.  It is an ordinary proof, with exact finite algebraic checks
and an independent review; it is not a new end-to-end Lean/Rocq certificate
or an exported derivation in an internally encoded ZFC calculus.

\subsection{The fixed integers and the complete quadratic system}
\label{int:literal}

Use the whole source34 \texttt{mixed-parity-unrounded-1064921} tuple.
Its ports are
\[
 \beta=2^{1064921}-1,\qquad K_0,\qquad M_C,\qquad M_F.
\]
The last port is the already shifted high mask.  All four literal integers
are supplied as lossless gzip payloads of minimal unsigned little-endian
bytes.  Their bit lengths are respectively
$1064921,1065043,1064921,1064922$; their combined uncompressed payload
is $532479$ bytes.  The included standard-library verifier reconstructs
every byte from the retained $943$-tile list and checks the original hashes.
No coefficient is selected by assuming consistency.  The relevant fixed
facts are
\begin{equation}\label{int:fixed-bounds}
 \beta\geq17,\qquad 0<M_C<\beta<M_F<2\beta,
 \qquad K_0\geq0,\qquad \beta\equiv1\pmod3.
\end{equation}

The thirty-one core unknowns are
\[
\begin{gathered}
J_3,F,\alpha,z,f,h,r,v,w_3,\tau,\eta,\zeta,y,\gamma,o,j,\\
N,E,P,c,a,X,D,C_2,L,F_2,V,g,I,B_0,\Delta.
\end{gathered}
\]
There are also three private triples $\mathbf u,\mathbf v,\mathbf e$,
for a total of forty unknowns.  In the JSON packet, $v$ is
\texttt{v\_scale}, $B_0$ is \texttt{Baux}, and the triples are
\texttt{FA\_pos1,...,3}, \texttt{FC\_pos1,...,3}, and
\texttt{eta\_zeta\_pos1,...,3}.  Every coordinate ranges over all integers.
Define only the following polynomial abbreviations:
\begin{align*}
 J_0&=3J_3+1,& W&=4(3w_3+1),&q&=\beta J_0+1,\\
 C&=q-F-\alpha,& k&=\eta+\zeta,\\
 \mathcal R&=(q-F)(N-q)-CL+(M_C+qM_F)J_0.
\end{align*}
Also set $\operatorname{Pos}(x_1,x_2,x_3)=1+x_1^2+x_2^2+x_3(x_3+1)$.
The exact main equation is $Q_{40}^{\mathrm{aff}}=\sum_{\nu=1}^{25}r_\nu^2=0$,
where the twenty-five residuals are the following quadratic polynomials:
\begin{align*}
 r_1&=L-q^2+1,                         &r_2&=N-q(L+1),\\
 r_3&=X-qW,                           &r_4&=a-E-(2v+1)N,\\
 r_5&=E(X+1)-Xa,                      &r_6&=c-\eta+(E-a)k,\\
 r_7&=P-k-E(c-\eta),                  &r_8&=\tau^2-P^2+Pk-1,\\
 r_9&=D-X-ac-\gamma(4a+3),            &r_{10}&=\Delta-a^2-4a-3,\\
 r_{11}&=C_2-c^2,                     &r_{12}&=D^2-C_2\Delta-1.
\end{align*}
\begin{align*}
 r_{13}&=k-1-hE-\mathcal R,
   &r_{14}&=(K_0+W)C-z(q-1)-F,\\
 r_{15}&=F_2-f^2,
   &r_{16}&=V+c-fo,\\
 r_{17}&=V-jc+k-1-hE,
   &r_{18}&=V^2+g-y^2,\\
 r_{19}&=I-(2r+1)C_2,
   &r_{20}&=B_0-I^2,\\
 r_{21}&=B_0-\Delta(F_2-1),
   &r_{22}&=y^2-1-B_0g,\\
 r_{23}&=F\alpha-\operatorname{Pos}(\mathbf u),
   &r_{24}&=FC-\operatorname{Pos}(\mathbf v),\\
 r_{25}&=\eta\zeta-\operatorname{Pos}(\mathbf e).
\end{align*}
No abbreviation adds an existential unknown.  Expanding these rows gives
exactly the literal packet \texttt{quartic40.json}.  Each row has degree
two.  Only $r_8$ contains $\tau$, and its coefficient of $\tau^2$ is one;
therefore $[\tau^4]Q_{40}^{\mathrm{aff}}=1$, proving exact degree four on the actual
fixed-coefficient slice.

\subsection{What changed from the parent}
\label{int:parent}

In source64's $41$-variable quartic, make the affine substitutions
\[
 J_{\rm old}=3J_3+1,\qquad
 w_{{\rm scale},\rm old}=3w_3+1,
\]
then delete the residual $X-2u^2$ and its private coordinate $u$.
The other twenty-five rows are unchanged under those substitutions.
Since $\beta\equiv1\pmod3$, the new rows imply
\begin{equation}\label{int:orientation-residue}
 q\equiv2\pmod3,\qquad W\equiv1\pmod3,
 \qquad X=qW\equiv2\pmod3.
\end{equation}
The work is to prove that this residue restriction can replace the deleted
positivity witness without introducing spurious zeros.

\subsection{Bounds before any sign is chosen}
\label{int:bounds}

At an integer zero every residual is zero.  The form
$\operatorname{Pos}$ is strictly positive on $\mathbb Z^3$ and represents
every positive integer.  To see surjectivity, apply Legendre's three-square
theorem to $4n-3$, which is $1$ modulo $4$ and cannot have the exceptional
form $4^a(8b+7)$; its representation has one odd and two even squares.

The two guards $F\alpha>0$ and $FC>0$ force $F,\alpha,C$ to have a common
nonzero sign.  Thus $q=F+\alpha+C$ has that sign and $|q|\geq3$.
Write $Q=|q|$.  Because $q=\beta J_0+1$, one has $J_0\neq0$ and
$Q\geq\beta-1\geq16$.  Set
\[
 R=((q-1)(q-F)+\alpha)(q^2-1)+(M_C+qM_F)J_0.
\]
The first two rows give $L=q^2-1$, $N=q^3$, and hence $\mathcal R=R$.
The sign-free packing estimate is
\begin{equation}\label{int:R-bounds}
       Q(Q^2-1)<R<Q^4-2.
\end{equation}
Here is a direct reminder of its proof.  If $q=Q>0$, put
$G=Q(Q-F)-C$.  Then
\[
 Q+2\leq G\leq Q^2-Q-1,
 \quad0<(M_C+QM_F)J_0<(2Q+1)(Q-1).
\]
If $q=-Q<0$, put $F_0=-F$, $C_0=-C$; now
$G=Q(Q-F_0)+C_0$ satisfies
\[
 2Q+1\leq G\leq Q^2-2,
 \quad0<(QM_F-M_C)(Q+1)/\beta<2Q(Q+1).
\]
Multiplying the displayed bounds for $G$ by $Q^2-1$ gives
\eqref{int:R-bounds} in both cases.  This uses no sign of $X,Y,k$.

Let $s=2v+1$, $i=2r+1$, and $Y=sq^3$.  Rows $r_4,r_5$ give
$a=E+Y$ and $E=XY$; thus $a=(X+1)Y$.
The nonzero integer $W=12w_3+4$ has $|W|\geq4$.  Since $s$ is odd,
writing $U=|X|$, $Z=|Y|$, and $M=|E|$ gives
\begin{equation}\label{int:UZ-bounds}
 U\geq4Q\geq64,\qquad Z\geq Q^3\geq4096,
 \qquad M=UZ\geq4Q^4>3R+1.
\end{equation}
The last guard implies that $\eta,\zeta$ have a common nonzero sign.
Put
\[
 \xi=\operatorname{sgn}X,\quad
 \sigma=\operatorname{sgn}k,\quad
 \delta=\operatorname{sgn}Y,\quad
 K=|k|,\quad \eta_0=|\eta|,\quad c_0=|c|.
\]
Then $K\geq2$, $0<\eta_0<K$, and
\begin{equation}\label{int:c-ratio}
 c_0=KZ+\delta\eta_0,\qquad
 K(Z-1)<c_0<K(Z+1),\qquad
 \operatorname{sgn}c=\sigma\delta.
\end{equation}
Set $A=a+2$, $H=4a+3$, and define the positive absolute Pell parameters
\[
 B=|A|=(U+\xi)Z+2\xi\delta,\qquad
 P_*=|2XY^2+1|=2UZ^2+\xi.
\]
Their signs in these expressions are valid because $(U-1)Z>2$.
They satisfy
\begin{equation}\label{int:BP-bounds}
 B\geq(4Q-1)Q^3-2>Q^4>R,\qquad P_*>B\geq2,\qquad c_0>K.
\end{equation}
For instance $P_*-B\geq2UZ^2-(U+1)Z-3>0$ in
\eqref{int:UZ-bounds}.  These inequalities hold in every sign branch.

\subsection{A sign-uniform Pell power-recovery argument}
\label{int:signed-power}

For $b\geq2$, use the standard Pell polynomials
\[
 \chi_b(t)+\psi_b(t)\sqrt{b^2-1}
       =(b+\sqrt{b^2-1})^t.
\]
The following proof retains all signs until the last congruence.

\paragraph{Classify the first and main norms.}
The source rows imply
\[
 \tau^2-XY^2(XY^2+1)k^2=1,\qquad
 D^2-(B^2-1)c^2=1.
\]
Consecutive-product Pell classification gives positive indices $n,p$ with
\begin{equation}\label{int:classify}
 |\tau|=\chi_{P_*}(n),\quad K=2\psi_{P_*}(n),\quad
 |D|=\chi_B(p),\quad c_0=\psi_B(p).
\end{equation}
There is no hidden positive-$X$ hypothesis: if $X<0$, put $m=UZ^2$;
the first coefficient is $m(m-1)=(m-1)m$ and its positive Pell parameter
is $2m-1=P_*$.  Since $P_*>B$ and $c_0>K$, monotonicity of $\psi$
implies $p>n$.  Otherwise
$c_0=\psi_B(p)\leq\psi_{P_*}(n)=K/2$, a contradiction.

\paragraph{Establish the ordinary-rank hypotheses.}
Modulo $M=|E|$, the parameter $P_*$ is congruent to $\xi$.
The recurrence for $\psi$ gives
$\psi_{P_*}(n)\equiv\xi^{n-1}n\pmod M$.
The index row $k=R+1+hE$ therefore gives
\begin{equation}\label{int:first-index}
 \sigma\xi^{n-1}2n\equiv R+1\pmod M.
\end{equation}
If the sign on the left is positive, $n\geq(R+1)/2$.
If it is negative, then $2n\geq M-R-1>2R$, so $n>R$.
Thus in both branches $p>n\geq(R+1)/2\geq24$.
The elementary bound $\psi_B(p)\geq(2B-1)^{p-1}$ gives
\begin{equation}\label{int:rank-hypotheses}
 c_0>B(B^2-1)^2,\qquad c_0>2p,\qquad c_0>2R.
\end{equation}
The first already follows for $p\geq6$; the last also follows from
$c_0\geq\psi_B(2)=2B>2R$.

The ordinary strong equation is
\[
 (ic^2)^2=(B^2-1)(f^2-1),\qquad i=2r+1\neq0.
\]
The inherited ordinary strong-rank lemma, now with its size hypothesis
proved in every sign branch, supplies an $m$ such that
\begin{equation}\label{int:strong-rank}
 |f|=\chi_B(m),\qquad p\mid m,\qquad c_0\mid m,\qquad m>2p.
\end{equation}
The exact source is the retained
\texttt{PELL\_RELAXED\_AUXILIARY\_PROOF.md}; its signed use and the
next step are detailed in the retained odd-scale and signed-domain notes.
Only $B,c_0,|f|,|i|$ occur in its hypotheses.

\paragraph{Recover the main rank before orienting signs.}
Let $\rho=|i|c_0^2$.  The auxiliary equations and exposed quotients give
\[
 \rho^2V^2-(\rho^2-1)y^2=1,\qquad
 V\equiv-c\pmod{|f|},\qquad V\equiv-R\pmod{c_0}.
\]
For an odd auxiliary index $\ell=2t+1$ and a sign $e_0$, Pell
classification gives $V=e_0 Q_t(\rho^2)$, where
\[
 Q_t(1-B^2)=(-1)^t\psi_B(2t+1),\qquad
 Q_t(0)=(-1)^t(2t+1).
\]
Here $\rho^2\equiv1-B^2\pmod{|f|}$ and
$\rho^2\equiv0\pmod{c_0}$.  Squaring the first quotient congruence
removes both $e_0$ and the sign of $c$.  The inherited signed Pell
step-down theorem, using $0<2p<m$, gives
$\ell\equiv\pm p\pmod m$; the second congruence gives
$\ell\equiv\pm R\pmod{c_0}$.
Since $c_0\mid m$, one obtains $R\equiv\pm p\pmod{c_0}$.
The intervals in \eqref{int:rank-hypotheses} imply
\begin{equation}\label{int:p-equals-R}
                         p=R.
\end{equation}
This is precisely the sign-insensitive auxiliary interface used by the
parent, with its prerequisites newly established also for $X<0$.

\paragraph{Unwrap the first index.}
Since $n<p=R$, the negative-sign branch of \eqref{int:first-index},
which required $n>R$, is impossible.  In the positive branch write
$2n=R+1+v_0M$, with $v_0\geq0$.  If $v_0\geq1$, then
$M>3R+1$ again implies $n>R$.  Therefore
\begin{equation}\label{int:unwrap}
 \sigma\xi^{n-1}=1,\qquad 2n=R+1,\qquad R\text{ is odd}.
\end{equation}

\paragraph{One ratio estimate handles all signs.}
Put $r_0=n-1=(R-1)/2\geq2$.  The elementary upper and lower Pell
growth bounds give
\[
 \frac{c_0}{K}
 =\frac{\psi_B(2r_0+1)}{2\psi_{P_*}(r_0+1)}
 \geq\frac12\left(\frac{(2B-1)^2}{2P_*}\right)^{r_0}.
\]
For $U\geq64$ and $Z\geq6$,
\begin{equation}\label{int:uniform-ratio}
               \frac{(2B-1)^2}{2P_*}>U-3.
\end{equation}
Indeed $B\geq(U-1)Z-2$ and $P_*\leq2UZ^2+1$, and
\begin{align*}
 &(2(U-1)Z-5)^2-(U-3)(4UZ^2+2)\\
 &\qquad=4(U+1)Z^2-20(U-1)Z+31-2U\\
 &\qquad=U(4Z^2-20Z-2)+(4Z^2+20Z+31)>0.
\end{align*}
Both final parentheses are positive for $Z\geq6$.
Together with \eqref{int:c-ratio}, this implies
\[
 Z+1>\frac{(U-3)^{r_0}}2\geq\frac{61^{r_0}}2.
\]
As $61^{r_0}/2-1>2^{2r_0+1}$ for $r_0\geq2$, we obtain
\begin{equation}\label{int:large-Y}
                            Z>2^R.
\end{equation}
The verifier checks the exact polynomial identity and positive coefficient
certificates after shifting $U=64+u$, $Z=6+z$; the displayed argument is
the universal proof rather than a finite-box test.

\paragraph{Exclude reciprocal roots and recover a signed power.}
Let $e=\operatorname{sgn}c$ and write $D=\lambda\chi_A(R)$ with
$\lambda\in\{1,-1\}$.  Since $R$ is odd,
$\psi_A(R)=\psi_B(R)$, even for $A<0$; hence $c=e\psi_A(R)$.
The polynomial Pell congruences, valid for either sign of $A,H$, are
\begin{equation}\label{int:root-congruences}
 \chi_A(R)-a\psi_A(R)\equiv2^R\pmod H,\qquad
 \chi_A(R)+a\psi_A(R)\equiv2^{-R}\pmod H.
\end{equation}
The odd modulus makes $2$ invertible.  At index one the two congruences
are $A-a=2$ and $2(A+a)\equiv1\pmod H$; the recurrence proves all indices.

If $\lambda=-e$, the row $D-ac=X+\gamma H$ gives
$X2^R\equiv-e\pmod H$.  But
\[
 1<U2^R<UZ<4(U-1)Z-4\leq|H|-1,
\]
by \eqref{int:large-Y}.  Neither sign of $X2^R$ can be congruent to
either unit modulo $|H|$.  This branch is impossible.
If $\lambda=e$, the same row gives $X\equiv e2^R\pmod H$; now
\[
 |X-e2^R|\leq U+2^R<U+Z<4(U-1)Z-3\leq|H|.
\]
This multiple of $H$ must be zero.  We have proved, before assuming
any positive sign of $X$,
\begin{equation}\label{int:signed-power-conclusion}
                      X=e2^R,\qquad R\text{ odd}.
\end{equation}

\subsection{Orientation, soundness, and canonical completeness}
\label{int:finish}

Because $R$ is odd, $2^R\equiv2\pmod3$, while
$-2^R\equiv1\pmod3$.  Combining
\eqref{int:signed-power-conclusion} with
\eqref{int:orientation-residue} forces $e=1$ and $X=2^R$.
The integer
\[
                         u=2^{(R-1)/2}
\]
therefore restores the deleted row $X=2u^2$.
Restore also $J_{\rm old}=3J_3+1$ and
$w_{{\rm scale},\rm old}=3w_3+1$.
Every parent residual now vanishes literally.  The source64 quartic41
theorem gives the finite ZFC refutation.  Its theorem is invoked only
after its square guard has been restored, so no positive-$X$ argument
has been used circularly.

Conversely, a finite refutation supplies the canonical source34 witnesses
used in quartic41 completeness.  With $d=1064921$, they satisfy
\[
 \begin{gathered}
 q=2^t,\quad t\text{ odd},\quad R\equiv3\pmod4,\quad
 X=2^R,\quad R-t\geq2\text{ even},\\
 J_{\rm old}=\frac{2^t-1}{2^d-1},\qquad
 w_{{\rm scale},\rm old}=2^{R-t-2}.
 \end{gathered}
\]
Modulo $3$ the first old coordinate is $1$, since $\beta\equiv1$
and $2^t\equiv2$.  The second is $1$ because its exponent is
nonnegative and even.  Thus
\[
 J_3=\frac{J_{\rm old}-1}{3},\qquad
 w_3=\frac{w_{{\rm scale},\rm old}-1}{3}
\]
are integers.  Retain all other coordinates and discard $u$ to obtain a
zero of $Q_{40}^{\mathrm{aff}}$.  It is the canonical completeness family that is
preserved; no claim is made that every parent zero lifts to the two
affine restrictions.  This proves Theorem~\ref{int:theorem} for $Q_{40}^{\mathrm{aff}}$.

\subsection{The smaller direct-coordinate quartic}
\label{int:direct}

For the preferred main quartic, retain the parent $J,w_{\rm scale}$ and
replace only $X$ by $3X_3-1$, while deleting $X-2u^2$ and $u$.
Denote this polynomial by $Q_{40}^{\mathrm{dir}}$; it is the complete
packet in \texttt{quartic40\_direct.json}.
Equivalently, use all twenty-five rows of Subsection~\ref{int:literal}
with the core list modified by $J_3\mapsto J$, $w_3\mapsto w_{\rm scale}$,
and $X\mapsto X_3$, and now use the abbreviations
\[
 J_0=J,\qquad q=\beta J+1,\qquad W=4w_{\rm scale},
 \qquad X=3X_3-1.
\]
The other abbreviations and the nine private guard coordinates are
unchanged.  The resulting sum of squares has $40$ unknowns, $25$
quadratic residuals, exact degree four, and $393$ expanded symbolic
monomials.  Its coefficient of $\tau^4$ remains one.

Here $X\equiv2\pmod3$ gives $X\neq0$ immediately.  The retained equation
$X=qW$ and $q\neq0$ then give $W\neq0$, so $|W|\geq4$ and
$|X|\geq4|q|$.  Thus every bound and every step of the sign-uniform
argument applies without change, recovering $X=e2^R$ with $R$ odd.
The residue selects $e=1$, after which
$u=2^{(R-1)/2}$ and $X_{\rm old}=3X_3-1$ restore the parent.
For completeness, canonical $X_{\rm old}=2^R$ with $R$ odd gives the
integer $X_3=(2^R+1)/3$; keep $J,w_{\rm scale}$ and all remaining
coordinates and discard $u$.  This direct version needs no separate
residue restriction on $J,w_{\rm scale}$ or on $\beta$.

The two-affine-coordinate version is retained for the projections below:
there $X=Wq$ is a monic polynomial definition.  Directly solving
$3X_3-1=Wq$ would require division by $3$, which is not an allowed
integer-polynomial restoration and is not used here.

\subsection{The exact sextic and octic projections}
\label{int:projections}

Deleting a row $v-P(\text{other coordinates})$ and substituting its
integer-polynomial right-hand side everywhere gives a bijection of integer
zero sets.  The inverse restores $v$.  The transformations below use the
centered-root and triangular elimination ideas of source62; their purpose
here is to transfer the new deletion to additional degree slices.

\paragraph{Thirty-four integers, degree six.}
Make the triangular change
$\tau_{\rm old}=\tau_{\rm tail}+E(c-\eta)$ and eliminate
\[
 \begin{gathered}
 P=\eta+\zeta+E(c-\eta),\quad N=q(L+1),\quad X=Wq,\\
 C_2=c^2,\qquad F_2=f^2,\qquad g=y^2-V^2.
 \end{gathered}
\]
The first norm becomes exactly
\[
 \tau_{\rm tail}^2+
 (2\tau_{\rm tail}-\eta-\zeta)E(c-\eta)-1.
\]
The leading quartic terms cancel, so every surviving row has degree at
most three.  Six coordinates and six defining rows disappear, leaving
thirty-four integers and nineteen rows.  The only row involving $r$ is
$I-(2r+1)c^2$, so the coefficient of $r^2c^4$ in their sum of squares
is exactly $4$.  The equation in \texttt{sextic34.json} has exact degree six.

\paragraph{Twenty-nine integers, degree eight.}
Put $k_0=1-V+jc+hE$ and eliminate the eleven coordinates
\[
\begin{gathered}
 L=q^2-1,\quad N=q^3,\quad X=Wq,\quad C_2=c^2,\quad F_2=f^2,\\
 g=y^2-V^2,\quad B_0=I^2,\quad
 D=Wq+ac+\gamma(4a+3),\\
 \Delta=a^2+4a+3,\quad
 \zeta=k_0-\eta,\quad P=k_0+E(c-\eta).
\end{gathered}
\]
Every right-hand side uses retained coordinates only.  The index-exposure
row $r_{17}$ supplies the monic definition of $\zeta$, while $r_7$
supplies the restoration of $P$.  Fourteen rows remain, all of degree at
most four.  The ordinary strong row
\[
 I^2-(a^2+4a+3)(f^2-1)
\]
is the only row containing both $a$ and $f$; its square contributes
$a^4f^4$ with coefficient $1$.  Therefore the equation in
\texttt{octic29.json} has exact degree eight.  Both projections preserve
the exact fixed tuple and inherit the same consistency equivalence.

\subsection{A general principle and what the comparison means}
\label{int:principle}

\begin{proposition}[Orientation after arithmetic recovery]
Suppose a polynomial system, without a designated sign guard, has an
arithmetic proof that every integer zero satisfies $X=\pm b^r$, where
$r$ lies in a specified set of integer indices.  Let $m$ be a fixed modulus
such that the possible residues of $b^r$ and $-b^r$ are disjoint.  If
affine restrictions on existing coordinates force $X$ into the positive
residue class, preserve a complete canonical witness family, and the
deleted guard can be restored from the recovered positive power, then
they replace that guard without introducing new unknowns.
\end{proposition}
\begin{proof}
Disjointness excludes the negative branch of the recovered power.
The original sign guard can then be restored if it follows from the
recovered positive power, as $X=2u^2$ does here for odd $r$.
Canonical completeness supplies the converse existence implication.
\end{proof}

In this construction $b=2$, $m=3$, and the index is odd: the two residue
sets are the singletons $\{2\}$ and $\{1\}$.  The important requirement
is the sign-uniform recovery theorem; simply dropping a square guard
and imposing a residue would be unjustified without it.

At the inspected source pin the comparable variable counts are
$41,35,30$ at degrees $4,6,8$.  The new counts are $40,34,29$,
with respectively $25,19,14$ residuals.  Their exact expanded sums of
squares, with the four fixed ports left as coefficient symbols, have
$393,614,1238$ monomials when the direct quartic is selected; the
two-affine-coordinate quartic used for projections has $614$.
The prior $41$-variable quartic has $388$.
The affine restrictions enlarge these expanded
expressions.  Thus the gain is in unrestricted-integer variables and
residual counts at fixed degree, not a decrease in every size measure.
Native coefficient payload size is unchanged; no new bound for nine
positive or ten unrestricted integer unknowns is asserted.

The standard-library builder and a separately authored verifier both check
all literal affine substitutions, all twenty-five restoration identities
for each projection, exact symbolic expansions, liveness, degrees,
nonzero degree certificates, source hashes, and the actual coefficient
payloads.  Both pass in normal and optimized Python execution.  These
checks establish the finite algebraic claims.  The unbounded Pell,
native-compiler, and syntactic-refutation implications remain the stated
mathematical proof and its named dependencies.

\section{Transferring a strong Pell root into a transport gap}
\label{sec:root-transfer}

This section gives a polynomial in twelve strictly positive integer
unknowns of exact degree $84$, with a literal $103$-operation arithmetic
circuit.  Its predecessor in package 61 has degree $118$ and $104$
operations.  A thirteen-positive-variable companion has degree $30$ and
$94$ operations, against package 61's degree $34$ and $93$ operations.
These comparisons concern the actual circuits, count all live binary
additions, subtractions and multiplications, and treat the complete fixed
coefficient numerals as leaves.

The new twelve-variable construction removes the normalized strong Pell
root through the first transport carrier.  A positive upper approximation
to that root makes both recovered root signs harmless.  Strengthening an
existing transport multiplier makes the inverse coordinate positive.
Thus the construction does not require a new index congruence for the
auxiliary Pell ordinate.  The smaller thirteen-variable companion combines
package 61's transport carrier with the adaptive modulus of package 65.

\subsection{Fixed coefficients and the precise inherited boundary}

Use the whole compact coefficient tuple of package 44:
\[
 \beta=2^{590069}-1,\qquad K_0,\qquad M_C,\qquad M_F.
\]
The four integers are delivered as lossless payloads in
\texttt{coefficients/}; $M_F$ is already the shifted source mask.
They are constants, not existential parameters.  The fixed numerical
binding and the unchanged source proof are recorded separately from the
new reductions.

We use package 61's normalized positive-source interface.  With positive
coordinates $J,F,\alpha,z,f,h,i,T_{\rm old},s,w,\tau,\eta,\zeta,y,
\gamma_{\rm old}$, define
\begin{align*}
 q&=\beta J+1,& C&=q-F-\alpha,\\
 X&=wq,&Y&=sq^3,\\
 E&=XY,& k&=\eta+\zeta,\\
 c&=kY+\eta,&a&=E+Y,\\
 H&=4a+3,&\Delta&=a^2+H=(a+2)^2-1,\\
 \ell&=EkY,&I&=ic^2.
\end{align*}
\[
 R=\big((q-1)(q-F)+\alpha\big)(q^2-1)+(M_C+qM_F)J.
\]
The six source equations are
\begin{align}
 \tau^2-\ell(\ell+k)&=1,&
 D^2-\Delta c^2&=1,\qquad D=X+ac+H\gamma_{\rm old},\label{eq:rt-source-a}\\
 f^2-\Delta I^2&=1,&
 \Delta(f^2-1)(V^2-y^2)+y^2&=1,\label{eq:rt-source-b}\\
 k-hE-R&=1,&
 (K_0+w)C+(q-F)-z(q-1)&=1,\label{eq:rt-source-c}
\end{align}
where $V=c(T_{\rm old}f-1)-Rf^2$.
The inherited compiler and K215 calibration identify existence of a
positive solution with a finite syntactic ZFC refutation.  This source
interface, rather than a finite experiment, is the semantic premise used
below.

Put $B=\beta+1$.  We use the source's canonical completeness promises
\[
 q=B^N,\quad J=1+B+\cdots+B^{N-1},\quad
 N=581197^e>1,\quad R\equiv3\pmod4,\quad X=2^R.
\]
The first Pell index is $n_0=(R+1)/2\ge3$.  A normalized strong solution
can be chosen together with a positive auxiliary solution of the form
\[
 V=f^2(cT-R)-c,\qquad y=V+u,\qquad T,u>0.
\]
This is the normalized source and canonical construction of package 61's report, Sections 2 and 6, together with its positive auxiliary reconstruction lemma (Section 3, source label \texttt{lem:auxcomplete});
no stronger congruence for $y$ is required here.

For arbitrary positive outer coordinates, once $C>0$ has been proved,
the authentic masks give
\begin{equation}
 q\ge16,\qquad 3q<R<q^4-2.\label{eq:rt-Rbound}
\end{equation}
These bounds do not presuppose $q\mid X$.

\subsection{A weighted transport carrier and the adaptive main carrier}

Initially supply $X,Y$ directly.  Let $j,F,\alpha,f,i,T,X,Y,r,\eta,
\zeta,u,\gamma$ be positive, put $J=3j+1$, and define the registers above.
Let $h_*>0$ be an integer polynomial register, and set
\begin{align}
 W&=1+(q-1)h_*,&M&=q^3(q-1),\nonumber\\
 B_t&=(qK_0+X)C-qF,&
 Z&=M(\ell+k+r)+(q-1)Y-q^3WB_t,\label{eq:rt-transport}\\
 A_0&=Z^2-M^2\bigl(\ell(\ell+k)+1\bigr).\nonumber
\end{align}
The adaptive main carrier is
\begin{align}
 S&=4qY,& U&=qS=4q^2Y,& r_0&=k-R-1,\nonumber\\
 D_*&=U(X+ac)+H(U\gamma+SX-r_0),\label{eq:rt-main}\\
 A_1&=D_*^2-U^2(1+\Delta c^2).\nonumber
\end{align}
Keep the other two residuals
\begin{align}
 A_s&=f^2-\Delta I^2-1,\nonumber\\
 V&=f^2(cT-R)-c,\nonumber\\
 A_a&=(V+u)^2-\Delta(f^2-1)u(2V+u)-1.\label{eq:rt-aux}
\end{align}

\paragraph{Weighted-carrier lemma.}
Every positive simultaneous zero of $A_0,A_1,A_a,A_s$, with any positive
integer $h_*$, reconstructs a positive solution of
\eqref{eq:rt-source-a}--\eqref{eq:rt-source-c}.  Conversely, canonical
source witnesses lift whenever the chosen multiplier satisfies $h_*\ge k$.

\begin{proof}
If a rational number has integral square, it is an integer.  Apply this
to $Z/M$: $A_0=0$ supplies an integer $\tau$ with
$\tau^2=\ell(\ell+k)+1$.  Since $M\mid Z$, reduction modulo $q^3$
gives $q^3\mid Y$.  Reduction modulo $q-1$ gives $(q-1)\mid B_t$,
because $W\equiv q^3\equiv1\pmod{q-1}$.  Define the proof-only integers
$s=Y/q^3>0$ and $t=B_t/(q-1)$.  Dividing the carrier gives
\[
 \tau=\ell+k+r+s-Wt.
\]
As $|\tau|<\ell+k$, we obtain $Wt>r+s>0$, hence $t>0$ and $B_t>0$.
If $C\le0$, then $B_t=(qK_0+X)C-qF<0$, a contradiction.  Thus $C>0$
and \eqref{eq:rt-Rbound} is now available.

Write $A=a+2$ and $P_0=2XY^2+1$.  The main numerator is positive before
either $q\mid X$ or rank recovery: since $r_0<k$, $c>kY$ and
$UaY>H$, we have $D_*>Uac-Hk>0$.  Thus $A_1=0$ gives a positive integer
$D=D_*/U$ with $D^2-\Delta c^2=1$.
Pell classification supplies
\[
 D=\chi_A(p),\quad c=\psi_A(p),\qquad
 |\tau|=\chi_{P_0}(n_0),\quad k=2\psi_{P_0}(n_0).
\]
Here $p\ge2$ and $n_0\ge2$.  Indeed $n_0=1$ would give $k=2$ and
$c=2Y+1<2A=\psi_A(2)$, which is impossible.  Consequently
\[
 k>8XY^2,\qquad c>8XY^3\ge8q^9>2R,\qquad n_0<p,\qquad c>2p.
\]
The normalized strong residual implies $f=\chi_A(m)$,
$I=\psi_A(m)$, $c\mid m$, and $m>2p$; this is the normalized-rank
lemma of package 61.  Put $S_*=\Delta I$.  The auxiliary residual gives
\[
 S_*^2V^2-(S_*^2-1)y^2=1,\qquad y=|V+u|>0,
\]
with $V\equiv-c\pmod f$ and $V\equiv-R\pmod c$.
The signed auxiliary-index lemma, with $0<p,R<c/2$, yields $p=R$.
The rank lemmas used here do not require the deleted scale or index
quotient, nor the sign of $V$.

It remains to extract the adaptive congruence without discarding a factor
of three.  Since $590069$ is odd, $\beta\equiv1\pmod3$ and
$J=3j+1$ imply $q\equiv2\pmod3$.  Because $q^3\mid Y$,
$\gcd(q,H)=1$.  Moreover $H=4Y(X+1)+3$ is odd and
$\gcd(H,Y)=\gcd(3,Y)$.  Hence
\[
 g:=\gcd(H,U)=\gcd(H,S)\in\{1,3\}.
\]
Integrality of $D_*/U$ gives $U/g\mid SX-r_0$, hence $S/g\mid r_0$.
The first Pell recurrence gives
\[
 k\equiv2n_0\pmod{4XY^2}.
\]
Since $S\mid4XY^2$ and $S/g\ge4q^4/3>R+1$,
the inequalities $1\le n_0<p=R$ force
$2n_0-R-1=0$.  Thus $r_0=k-2n_0$ is a positive multiple of $4XY^2$.
Also $U\mid4XY^2$, since $q^2\mid Y$; therefore $U\mid r_0$.
It follows from $U/g\mid SX$ that $q\mid gX$.
Because $g$ divides $3$ and $3\nmid q$, this gives $q\mid X$.

Now $w=X/q>0$, $h=r_0/(XY)>0$, and
$\gamma_{\rm old}=\gamma+w-r_0/U$ are integers.  Since
$D>(a+1)c$ and $c>X$, the identity
$H\gamma_{\rm old}=D-ac-X$ proves $\gamma_{\rm old}>0$.
Since $B_t=q((K_0+w)C-F)>0$ and $(q-1)\mid B_t$, the integer
$z_{\mathrm{new}}=((K_0+w)C-F)/(q-1)$ is positive; set $z=z_{\mathrm{new}}+1$.
Finally replace $\tau$ by $|\tau|$, put $T_{\rm old}=fT$, and retain
$y=|V+u|$.  These are the required source coordinates.

For completeness, use canonical source witnesses.  Their odd repunit
exponent gives $J\equiv1\pmod3$, so $j=(J-1)/3>0$ is integral.
The source transport coordinate has $z\ge2$, hence
$t=q(z-1)\ge q$.  Put
\[
 r=Wt+\tau-\ell-k-s.
\]
Since $\tau>\ell$, $s<k$, and $h_*\ge k$, this gives
$r>qW-2k>0$.  It makes $A_0=0$.
The canonical first index gives $r_0$ a positive multiple of $4XY^2$;
take $\gamma=\gamma_{\rm old}-w+r_0/U$.
It is positive because
$r_0/U\ge XY/q^2=wY/q>w$.  Its substitution makes $A_1=0$.
The positive auxiliary reconstruction supplies $T,u$, and the two
remaining residuals vanish.
\end{proof}

Taking $h_*=k$ defines the thirteen-variable companion
\[
 P_{13}=A_0^2+A_1^2+A_a^2+A_s^2.
\]
Its residual degrees are respectively $15,13,13,14$, so its exact total
degree is $30$.  The factor $4$ in $S=4qY$ is used in the proof:
removing it without another argument loses the bound $S/g>R+1$.

\subsection{The twelve-variable polynomial}

Supply exactly
\[
 j,F,\alpha,i,T,X,Y,n,\eta,\zeta,u,\gamma\in\mathbb Z_{>0}.
\]
Use the same common registers, but now define
\begin{align}
 D_f&=1+\Delta I^2,& L&=(a+2)I,& h_*&=ikc,\nonumber\\
 W&=1+(q-1)ikc,&
 Z_0&=M(\ell+k+n)+(q-1)Y-q^3WB_t,\label{eq:rt-newregs}\\
 \overline Z&=Z_0+ML,& D_\tau&=\ell(\ell+k)+1.\nonumber
\end{align}
The first residual is the ordinary polynomial
\begin{equation}
 N=\bigl(\overline Z^2-M^2(D_f+D_\tau)\bigr)^2
                  -4M^4D_fD_\tau.\label{eq:rt-N}
\end{equation}
Retain $A_1$ from \eqref{eq:rt-main} and replace the auxiliary registers by
\begin{equation}
 V=D_f(cT-R)-c,\qquad Q=\Delta(D_f-1),\qquad
 A_a=(V+u)^2-Qu(2V+u)-1.\label{eq:rt-newaux}
\end{equation}
The recommended equation is
\begin{equation}
 \boxed{P_{12}=N^2+A_1^2+A_a^2=0.}\label{eq:rt-P12}
\end{equation}
Every named register is a polynomial abbreviation.  There are twelve,
not twelve plus the number of registers, quantified unknowns.

\paragraph{Root-transfer theorem.}
The polynomial \eqref{eq:rt-P12} has a strictly positive zero if and only
if the inherited source has a strictly positive zero.  Consequently its
lack of positive zeros is equivalent in ZFC to $\operatorname{Con}(\mathrm{ZFC})$,
subject exactly to the inherited compiler/calibration interface above.

\begin{proof}
The inequalities needed to recover roots hold on every positive input,
before any residual is assumed zero.  Since $\ell<ac$ and $k<c$,
\[
 \ell(\ell+k)<a(a+1)c^2<\Delta i^2c^4,
\]
and therefore $D_f>D_\tau>0$.  Also
\[
 L^2-D_f=I^2-1>0,
\]
because $c\ge3$ and $I=ic^2>1$.

At a zero of $N$, $\overline Z=0$ is impossible: it would imply
$M^4(D_f-D_\tau)^2=0$.  The identity
\begin{align*}
 N={}&\bigl(\overline Z^2+M^2(D_f-D_\tau)\bigr)^2
                     -4M^2\overline Z^2D_f
\end{align*}
then shows that the rational number
\[
 f_*=\frac{\overline Z^2+M^2(D_f-D_\tau)}{2M\overline Z}
\]
has square $D_f$.  It is a nonzero integer by rational-square
integrality.  Furthermore
\[
 \left(\frac{\overline Z-Mf_*}{M}\right)^2=D_\tau,
\]
so this second rational root is also an integer.
Set $f=|f_*|$ and
\[
 r=n+L-f_*.
\]
The strict inequality $L>|f_*|$ gives $r>0$ for both signs of $f_*$.
The first carrier becomes $Z=\overline Z-Mf_*$, so $A_0=0$.
The strong residual is zero because $f^2=D_f$.
All other registers involving the strong root use its square only;
thus the supplied $A_1=A_a=0$ give the remaining weighted-carrier
equations.  The lemma reconstructs a positive source solution.

Conversely, use canonical source witnesses and their positive strong root
$f$.  Choose $T,u$ by positive auxiliary reconstruction, and $r,\gamma$
by the weighted-carrier lemma with $h_*=ikc$.  This multiplier is at least
$k$.  The upper approximation gap satisfies
\[
 0<L-f=\frac{I^2-1}{L+f}<\frac{I}{a}
       =\frac{ic^2}{a}<ikc,
\]
because $c=kY+\eta<k(Y+1)\le ka$ for $X,Y\ge1$.
The transport construction gives
\[
 r>qW-2k>ikc.
\]
The last inequality follows directly from
$q\ge16$ and $ic\ge3$.
Hence $n=r-L+f$ is a strictly positive integer.
Its substitution gives $\overline Z=M(\tau+f)$, and the two source
squares imply $N=0$.  The other two residuals already vanish.
This proves completeness, including positivity of the new coordinate.
\end{proof}

The theorem proves an equivalence and does not decide either side.
The norm is used with its zero-denominator exclusion and both signed-root
branches; an abstract appeal to independent square classes would not
supply these facts.

\subsection{Exact degree, arithmetic cost, and finite certificates}

Write
\[
 q_0=3\beta j,\quad k_0=\eta+\zeta,\quad
 B_{t,0}=(K_0q_0+X)(q_0-F-\alpha)-q_0F.
\]
For the twelve-variable equation,
\[
 \deg M=4,\quad \deg L=7,\quad \deg D_f=14,\quad
 \deg D_\tau=8,\quad \deg Z_0=10.
\]
The cancellation-resolved middle expression in \eqref{eq:rt-N} is
\begin{equation}
 \overline Z^2-M^2(D_f+D_\tau)
 =Z_0^2+2MLZ_0-M^2(1-I^2+D_\tau).\label{eq:rt-centeridentity}
\end{equation}
It has degree $21$; only its middle summand reaches that degree.
The subtracted norm tail $4M^4D_fD_\tau$ has degree $38$.  Thus $N$ has
exact degree $42$, with leading homogeneous form
\[
 \operatorname{in}_{42}(N)
 =4q_0^{16}i^4k_0^8X^2Y^8B_{t,0}^{\,2}.
\]
The other residual degrees are $13$ and $37$, respectively.  Consequently
\begin{equation}
 \operatorname{in}_{84}(P_{12})
 =16q_0^{32}i^8k_0^{16}X^4Y^{16}B_{t,0}^{\,4}.
 \label{eq:rt-leading84}
\end{equation}
In particular the coefficient of
$j^{32}i^8\eta^{16}X^8Y^{16}\alpha^4$ is exactly
$16(3\beta)^{32}>0$.  Choosing this minimum $j$ exponent forces the
$-X\alpha$ term from each copy of $B_{t,0}$, so the other fixed ports
cannot cancel it.  This proves exact degree after the actual coefficient
specialization.

For $P_{13}$, the unique highest residual degree is $15$ and
\[
 \operatorname{in}_{15}(A_0)=-2q_0^8k_0^2XY^2B_{t,0},\qquad
 \operatorname{in}_{30}(P_{13})=4q_0^{16}k_0^4X^2Y^4B_{t,0}^{\,2}.
\]
Its coefficient of $j^{16}\eta^4X^4Y^4\alpha^2$ is
$4(3\beta)^{16}>0$.

The delivered twelve-variable evaluator uses $60$ multiplications,
$27$ additions and $16$ subtractions, totaling $103$ live operations.
The three-equation version uses $98$ operations and residual degrees
$42,13,37$.  The thirteen-variable evaluator uses $94$ operations,
or $87$ for its four-equation system.
The JSON graphs contain no divisions or variable exponents.

The verifier checks literal graph liveness, every fixed-payload digest,
byte-identical circuit regeneration, exact prefix identities,
independent cut-ring norm and finalizer identities, and the nonzero
leading forms above.  These are finite certificates of algebra and data.
The unbounded existence equivalence follows from the written proofs
and the inherited source theorem, not from bounded tests.  No new
end-to-end Lean/Rocq certificate or exported object-language ZFC
derivation is claimed for these literal polynomials.

\section{Three-root transport: eleven positive variables, degree 325}
\label{rt11:section}

There is a division-free polynomial in eleven strictly positive integer
unknowns of exact total degree $325$, evaluated by $113$ binary arithmetic
operations, whose lack of positive solutions is equivalent in ZFC to
$\operatorname{Con}(\mathrm{ZFC})$. At the pinned repository state the
eleven-positive equation in package~52 has degree $954$ and the same
$113$-operation cost. The construction below removes $629$ degrees while
retaining that operation count. Its two residuals have degrees $312$ and
$13$ and together require $110$ operations.

The new ingredient is a particular combination of the normalized positive
auxiliary reconstruction, the weighted first carrier, and an elementary
denominator argument. Conjugate norms, positive majorants, radical
separation, quotient absorption and unit-product finalizers already occur
in package~52 and its predecessors; those general methods are not claimed
as new. In package~52 the auxiliary radicand depends on the signed strong
Pell root and the main Pell root is eliminated inside the same norm.
Here the normalized auxiliary expression depends only on its square.
The first, strong and auxiliary roots therefore enter a direct octic
sum norm, while the inexpensive adaptive main equation remains separate.

\subsection{Fixed data and the literal equation}

Use precisely the coherent tuple $(\beta,K_0,M_C,M_F)$ of
Section~\ref{sec:root-transfer}, with $\beta=2^{590069}-1$ and the already
shifted $M_F$. These are fixed integer numerals, not parameters.
Their payloads and source hashes are supplied again in this component.
The authentic mask inequalities used below are
\[
 0<M_C<\beta,\qquad 0<M_F<2\beta,\qquad K_0>0.
\]
The eleven unknowns are exactly
\[
 j,F,\alpha,i,T,X,Y,n,\eta,\zeta,\gamma\in\mathbb Z_{>0}.
\]
Define polynomial registers
\begin{align*}
 J&=3j+1,&q&=\beta J+1,&C&=q-F-\alpha,\\
 E&=XY,&k&=\eta+\zeta,&c&=kY+\eta,\\
 a&=E+Y,&H&=4a+3,&\Delta&=a^2+H=(a+2)^2-1,\\
 \ell&=EkY,&M&=q^3(q-1),&B_t&=(qK_0+X)C-qF,\\
 D_\tau&=\ell(\ell+k)+1.&&&
\end{align*}
Set
\begin{align}
 R&=\bigl((q-1)(q-F)+\alpha\bigr)(q^2-1)+(M_C+qM_F)J,
       \label{rt11:R}\\
 \widehat R&=q^3(q+F+\alpha),\nonumber
\end{align}
and then
\begin{align}
 I&=iq(q-1)c^2,&D_f&=1+\Delta I^2,&Q&=\Delta(D_f-1),
       \label{rt11:strong}\\
 \delta&=Q-1,&V&=D_f(cT-R)-c,&D_y^{\rm num}&=QV^2-1,
       \nonumber\\
 \overline V&=D_f(cT+\widehat R),&L&=3\overline V,&W&=1+(q-1)L.
       \label{rt11:majorant}
\end{align}
The first transport register is
\begin{equation}
 \overline Z=M(\ell+k+n+L)+(q-1)Y-q^3WB_t.
 \label{rt11:Z}
\end{equation}
The adaptive main carrier is unchanged:
\begin{align}
 r_0&=k-R-1,& S&=4qY,& U&=qS,\nonumber\\
 D_*&=U(X+ac)+H(U\gamma+SX-r_0),\nonumber\\
 A_1&=D_*^2-U^2(1+\Delta c^2).
 \label{rt11:main}
\end{align}

The following schedule evaluates the octic norm using only integer
polynomial operations:
\begin{align}
 \mathcal A&=M^2D_\tau,&
 \mathcal B&=M^2D_f,&
 \mathcal C&=M^2D_y^{\rm num},\nonumber\\
 T_n&=\delta^2\mathcal A\mathcal B,&
 R_n&=\delta\mathcal C,&
 J_n&=\delta(\overline Z^2-\mathcal A-\mathcal B)-\mathcal C,
       \nonumber\\
 N_*&=J_n^2-4\bigl(T_n+R_n(\mathcal A+\mathcal B)\bigr),
       \nonumber\\
 N_8&=N_*^2-64T_nR_n\overline Z^2.
 \label{rt11:norm}
\end{align}
The equation is
\begin{equation}
 \boxed{\ P_{11}=N_8+A_1+N_8A_1=0.\ }
 \label{rt11:equation}
\end{equation}
All the displayed registers are abbreviations, with no additional
quantified witnesses. The supplied JSON is a literal binary
$+,-,\times$ circuit; even multiplication by a small integer is charged.
The rational quantity
\[
 D_y=\frac{D_y^{\rm num}}\delta
\]
is used only in the proof and is never evaluated by the polynomial circuit.

\subsection{An elementary rational-root extraction lemma}

\begin{lemma}
\label{rt11:rational}
Let $d_1,d_2,d_3$ be distinct positive rational numbers. Suppose signed
real roots $x_s^2=d_s$ have a nonzero rational sum $t=x_1+x_2+x_3$.
Then all three $x_s$ are rational.
\end{lemma}
\begin{proof}
Put
\[
 u=\frac{t^2-d_1-d_2-d_3}{2},\qquad
 v=\frac{u^2-d_1d_2-d_1d_3-d_2d_3}{2t}.
\]
These are rational. Expansion gives
$u=x_1x_2+x_1x_3+x_2x_3$ and $v=x_1x_2x_3$. Consequently
\[
 d_1t+v=x_1(x_1+x_2)(x_1+x_3),\qquad
 x_1=\frac{d_1(d_1+u)}{d_1t+v}.
\]
The denominator is nonzero. Indeed, $x_1\ne0$, while either
$x_1=-x_2$ or $x_1=-x_3$ would contradict distinctness of the radicands.
The two cyclic formulas recover $x_2,x_3$ as rational numbers as well.
This proof accommodates every coincidence or dependence among square
classes; it assumes no independent specialized Galois automorphisms.
\end{proof}

\subsection{Unconditional radical separation and the positive guard}

All bounds in this subsection hold at every positive input tuple, before
assuming any residual vanishes. In particular they do not assume $C>0$,
$q^3\mid Y$, or a recovered Pell index. We have
\[
 q\ge16,\quad k\ge2,\quad c\ge3,\quad a\ge2,\quad
 \Delta\ge15,\quad I\ge c^2,\quad \ell<ac,\quad k<c.
\]
Thus
\[
 \ell(\ell+k)<a(a+1)c^2<\Delta I^2,
 \qquad 1<D_\tau<D_f<Q.
\]
Write $t=\sqrt{D_\tau}$ and $f=\sqrt{D_f}$ for positive real roots.
Then
\begin{equation}
 t<\ell+k<(a+1)c,
 \qquad f>(a+1)I\ge(a+1)c^2\ge3c^2.
 \label{rt11:tf}
\end{equation}

The integer $b=cT-R$ determines two cases. If $b=0$, then $V=-c$.
If $b\ne0$, then $|V|\ge D_f-c>1$. Hence $|V|>1$ in both cases and
\begin{equation}
 D_y=V^2+\frac{V^2-1}{Q-1}>V^2>1.
 \label{rt11:Dy}
\end{equation}
The radicands $D_\tau,D_f,D_y$ are pairwise distinct. For if
$D_y=D$ with $D\in\{D_\tau,D_f\}$, then
\[
 Q(D-V^2)=D-1,
\]
which is impossible because $0<D-1<Q$.

Let $y=\sqrt{D_y}>0$. If $b=0$, then $Q>c^2$ implies
$D_y<c^2+1$ and $y<c+1$, whereas \eqref{rt11:tf} gives
\[
 f-t>(a+1)c(c-1)>c+1.
\]
Thus $f>t+y$. If $b\ne0$, the bound $f>3c^2>c+2$ implies
$f^2-c>2f$, so
\[
 y>|V|\ge D_f-c>2f>f+t.
\]
In either case one root exceeds the sum of the other two. Therefore no
signed sum $\pm t\pm f\pm y$ is zero.

It remains to check that the polynomial $L$ bounds the two roots to be
transferred out of the positive coordinate. The authentic masks give
\[
 (M_C+qM_F)J<(q-1)(1+2q)<2q^2.
\]
The triangle inequality applied to \eqref{rt11:R} then gives
\[
 |R|<q^4+Fq^3+(\alpha+2)q^2
     <q^3(q+F+\alpha)=\widehat R,
\]
because $\alpha(q-1)>2$. Both sides are integers, and $D_f>c$; hence
\[
 |V|\le D_f(cT+|R|)+c
       <D_f(cT+\widehat R)=\overline V.
\]
Also $\overline V>D_f>f$ and, by \eqref{rt11:Dy},
$y<2|V|<2\overline V$. We obtain the uniform bounds
\begin{equation}
 L=3\overline V>f+y,
 \qquad L>c>k.
 \label{rt11:Lbound}
\end{equation}

\subsection{What the cleared norm forces}

The universal octic identity is
\begin{align}
 F_8(x;a,b,c)
 &:=\bigl((x^2-a-b-c)^2-4(ab+ac+bc)\bigr)^2-64abcx^2
       \nonumber\\
 &=\prod_{\epsilon_1,\epsilon_2,\epsilon_3\in\{-1,1\}}
       \bigl(x-\epsilon_1\sqrt a-\epsilon_2\sqrt b
                    -\epsilon_3\sqrt c\bigr).
 \label{rt11:genericnorm}
\end{align}
This is a polynomial identity, verified by multiplying the eight formal
conjugates. The center in \eqref{rt11:norm} also equals
\[
 \delta^2\bigl((\overline Z^2-\mathcal A-\mathcal B)^2
                       -4\mathcal A\mathcal B\bigr)
 -2\delta\mathcal C(\overline Z^2+\mathcal A+\mathcal B)
 +\mathcal C^2.
\]
Consequently
\begin{equation}
 N_8=\delta^4F_8
       (\overline Z;M^2D_\tau,M^2D_f,M^2D_y).
 \label{rt11:clearing}
\end{equation}
All three specialized radicands are positive, and $M,\delta>0$.
If $N_8=0$, one of its eight real factors vanishes. There are signed
roots satisfying
\begin{equation}
 \frac{\overline Z}{M}=\tau_*+f_*+y_*,
 \qquad \tau_*^2=D_\tau,\quad f_*^2=D_f,\quad y_*^2=D_y.
 \label{rt11:signedsum}
\end{equation}
The preceding separation proves that this rational sum is nonzero and
its three radicands are distinct. Lemma~\ref{rt11:rational} makes all
three signed roots rational. Since $D_\tau,D_f$ are integers,
rational-square integrality yields $\tau_*,f_*\in\mathbb Z$.

The auxiliary root becomes integral for a different, elementary reason.
Let $d$ be its positive reduced denominator. Equation
\eqref{rt11:signedsum} gives $d\mid M$, and
$y_*^2=D_y^{\rm num}/\delta$ gives $d^2\mid\delta$.
Every prime divisor of $M=q^3(q-1)$ divides $q(q-1)$ and therefore divides
$I=iq(q-1)c^2$ and $Q=\Delta^2I^2$. It cannot divide $\delta=Q-1$.
Thus
\[
 \gcd(M,\delta)=1,\qquad d=1,
\]
and $y_*$ is an integer. This uses only the prime support of $M$; no
claim that $M\mid Q$ is needed.

\subsection{Soundness and completeness}

\begin{theorem}
\label{rt11:theorem}
The equation \eqref{rt11:equation} has a strictly positive solution if
and only if the inherited normalized compiler has a strictly positive
solution. Hence, using the same compiler and refutation calibration,
ZFC proves that \eqref{rt11:equation} has no positive solution if and
only if $\operatorname{Con}(\mathrm{ZFC})$.
\end{theorem}
\begin{proof}
First, the single equation is equivalent to $N_8=A_1=0$. Indeed,
\[
 P_{11}=(1+N_8)(1+A_1)-1.
\]
Both residuals are squares modulo $4$: this follows directly from
\eqref{rt11:norm} and from $4\mid U$ in \eqref{rt11:main}.
At a zero of $P_{11}$, the two integer factors are therefore both $1$;
the alternative of both being $-1$ would force a residual to be
$-2\equiv2\pmod4$. The converse is immediate.

For soundness, recover the signed integral roots in
\eqref{rt11:signedsum}, and set
\[
 i_{\rm old}=iq(q-1),\quad f_{\rm old}=|f_*|,\quad
 y_{\rm old}=|y_*|,\quad u=y_{\rm old}-V,\quad
 r=n+L-f_*-y_*.
\]
Equation \eqref{rt11:Dy} gives $y_{\rm old}>|V|$, so $u>0$.
The uniform bound \eqref{rt11:Lbound} gives $r>0$ for every sign choice.
The strong equation is
\[
 f_{\rm old}^2-\Delta(i_{\rm old}c^2)^2=1.
\]
The auxiliary root equation is exactly
\[
 y_{\rm old}^2-Q(y_{\rm old}^2-V^2)=1,
\]
which is the normalized auxiliary residual with $y_{\rm old}=V+u$.
Finally the first transport register with weight $h_*=L$ becomes
\[
 Z=M(\ell+k+r)+(q-1)Y-q^3WB_t
   =\overline Z-M(f_*+y_*)=M\tau_*.
\]
It satisfies $Z^2=M^2D_\tau$. The weighted transport and adaptive main
carrier theorem from Section~\ref{sec:root-transfer} now applies, with
all its witnesses positive, and recovers a normalized compiler solution.
The signed roots must be retained in the definition of $r$; absolute
values are taken only in the strong and auxiliary witnesses.

For completeness, start with a canonical normalized compiler solution.
Keep its outer quantities, including $q,X,Y,k,c,a,R$ and its first
positive Pell root $\tau$. The canonical facts needed from the weighted
transport theorem are
\[
 R\equiv3\pmod4,\quad c=\psi_{a+2}(R),\quad
 t_0=\frac{B_t}{q-1}=q(z_{\rm source}-1)\ge q,
 \quad s=\frac{Y}{q^3}<k,\quad \tau>\ell.
\]
We must choose the strong Pell solution with the additional divisibility
$q(q-1)c^2\mid I$; this is not an assumption about a preselected source
witness. Put $m_0=q(q-1)c^2$ and consider
\[
 U_0=\begin{pmatrix}a+2&\Delta\\1&a+2\end{pmatrix}.
\]
Its determinant is $1$, so its reduction belongs to the finite group
$\mathrm{SL}_2(\mathbb Z/m_0\mathbb Z)$. A positive power is the
identity. Since the top-left and lower-left entries of $U_0^v$ are
$\chi_{a+2}(v)$ and $\psi_{a+2}(v)$, a positive multiple of that order
gives positive integers $f,I$ with
\[
 f^2-\Delta I^2=1,\qquad m_0\mid I.
\]
Set $i=I/m_0>0$. The positive auxiliary reconstruction lemma in the
pinned package~61 report (its label \texttt{lem:auxcomplete}) applies to
every such strong solution: its hypotheses are $R\equiv3\pmod4$,
$c=\psi_{a+2}(R)$ and $c\mid I$. It supplies positive $T,u$ and
$y=V+u>0$ satisfying the normalized auxiliary equation. Thus the stronger
divisibility preserves the required reconstruction promises.

Only after choosing these roots and $T$ do we compute
$\widehat R,\overline V,L,W$. Choose
\[
 r=Wt_0+\tau-\ell-k-s.
\]
Then
\[
 r>qW-2k=q+q(q-1)L-2k>L,
\]
using $q\ge16$ and $L>k$. Consequently
\[
 n=r-L+f+y>0.
\]
It follows that $\overline Z=M(\tau+f+y)$, so the norm vanishes.
The unchanged adaptive main construction supplies positive $\gamma$
with $A_1=0$. Its canonical repunit choice also gives $J=3j+1$ with
$j>0$, as in Section~\ref{sec:root-transfer}. This completes the lift.
The order of choices was the strong pair, the auxiliary pair, the
majorant, and then the transport coordinate; there is no circular
choice of witnesses or unproved enlargement of a special index.
\end{proof}

\subsection{Exact degree, circuit cost and verification boundary}

The ordinary total degrees in the eleven supplied unknowns are
\[
\begin{array}{c|rrrrrrrrr}
\text{register}&I&D_f&Q&V&L&W&\overline Z&D_y^{\rm num}&\mathcal C\\\hline
\text{degree}&7&18&22&22&22&23&28&66&74
\end{array}
\]
and $\deg\mathcal A=16$, $\deg\mathcal B=26$.
In the expanded expression for $N_*$, the term
$\delta^2\overline Z^4$ has degree $156$, while the other two groups
have degrees at most $152$ and $148$. The norm tail has degree $238$,
below $2\cdot156$. Thus the top part of $N_8$ is
$Q_0^4\overline Z_0^8$, of degree $312$.

For an explicit certificate, set
\[
 q_0=3\beta j,\qquad k_0=\eta+\zeta,\qquad
 B_{t,0}=(K_0q_0+X)(q_0-F-\alpha)-q_0F.
\]
The relevant leading forms are
\begin{align*}
 Q_0&=i^2q_0^4k_0^4X^4Y^8,\\
 \overline Z_0&=-3i^2q_0^{11}k_0^4X^2Y^6
                      (q_0+F+\alpha)B_{t,0},\\
 (A_1)_{13}&=32q_0^2k_0X^2Y^4
                   \bigl(4q_0^2Y\gamma+q_0^3(q_0-F)\bigr).
\end{align*}
The apparent degree-$14$ terms of $A_1$ cancel; the last expression is
verified by exact sparse polynomial expansion, not by gate-degree
propagation. Therefore the degree-$325$ part of $P_{11}$ is
\begin{align*}
 2^5 3^8\,i^{24}q_0^{106}k_0^{49}X^{34}Y^{84}
 (q_0+F+\alpha)^8B_{t,0}^8
 \bigl(4q_0^2Y\gamma+q_0^3(q_0-F)\bigr).
\end{align*}
Its coefficient of
\begin{equation}
 j^{108}i^{24}\eta^{49}X^{42}Y^{85}\alpha^{16}\gamma
 \label{rt11:coefficientmonomial}
\end{equation}
is exactly
\[
 2^7 3^8(3\beta)^{108}>0.
\]
Indeed the $\gamma$ term fixes the final factor, and minimum $j$ degree
reduces the two remaining parenthesized powers to
$X^8(F+\alpha)^{16}$, whose $\alpha^{16}$ coefficient is $1$.
This certifies degree $325$ on the actual fixed coefficient slice; no
specialization of $K_0,M_C,M_F$ can cancel the selected coefficient.

The literal circuit has $67$ multiplications, $29$ additions and
$17$ subtractions, for $113$ operations. Removing the three-operation
finalizer leaves $110$ operations for the two-residual system. The
normal and optimized-Python replay receipts check byte-identical circuit
regeneration, input and gate liveness, every polynomial register contract,
the formal octic identity, denominator clearing, rational-extraction
identities, the main-Pell cancellation, the exact degree certificate,
and the complete fixed payloads and retained source hashes.
An independent audit uses a separate sparse-ring implementation and
does not import the construction builder.

These checks certify finite algebra and data integrity. The existence
reduction is the proof above together with the preceding weighted carrier
theorem; the unbounded compiler and K215 refutation calibration remain
inherited mathematical dependencies. No end-to-end proof-assistant
certificate, global minimality theorem, or proof of consistency is claimed.


\section{Logical status, exact data, and reproducibility}
\label{sec:evidence}

\subsection{The proof chain and its limits}

There are three logically different kinds of evidence in this release.
First are ordinary unbounded mathematical proofs: the signed power
recovery and residue orientation, the root-transfer lemma, adaptive
congruence recovery, and soundness and completeness of the new
equations.  Second are inherited mathematical theorems: the Pell
rank/auxiliary interfaces, the native finite-window compiler, its fixed
coefficient adapters, and the K215 equivalence with the specified
syntactic refutation predicate.  Third are finite checks of literal
polynomials and data.  None of the latter replace the first two.

For $Q_{40}^{\mathrm{dir}}$ the proof restores the deleted square coordinate and hence
a zero of the parent $Q_{41}$.  This deliberately keeps the native
semantic interface unchanged.  The companion eliminations are exact
integer-polynomial bijections of zero sets.  The affine restrictions
used for completeness have a weaker and sufficient property: the
canonical witness family lies in their image.  We do not claim that
every parent zero lies in that image.

For $P_{12}$ the norm recovers a signed integer root $f_*$ and a
separate integer first root.  The gap uses $r=n+L-f_*$; the positive
strong coordinate is $|f_*|$.  These two roles must not be conflated.
The normalized source interface is then reconstructed with positive
scale, transport, and main-root quotients.  No division is used in the
literal equation, and every division appearing in the proof is
justified before its result is used as a witness.

For $P_{11}$ the analogous recovery uses three signed roots.  Pairwise
distinct rational radicands and a nonzero rational sum force all three
roots to be rational.  The first two are integral because their squares
are integers.  For the third root, its reduced denominator divides
$M=q^3(q-1)$ while the square of that denominator divides $Q-1$; the
engineered condition $\gcd(M,Q-1)=1$ removes it.  The proof restores
positive gap and auxiliary coordinates for every recovered sign choice.
Canonical completeness follows from a finite Pell-matrix period, so the
stronger divisibility is proved rather than presumed.

The active repository separately reports checked Lean consistency
equations and ongoing formalization of smaller arithmetic circuits.
Its README explicitly distinguishes a Lean theorem about a syntactic
predicate from an exported proof inside its object-language ZFC
calculus~\cite{active}.  This release does not run a Lean or Rocq kernel,
does not claim a new complete formal certificate, and does not transfer
the formal status of an existing equation to a different literal
polynomial.  The asserted equivalences are ordinary mathematical
proofs formalizable in ZFC with the named inherited results.  The
independent reviews here are internal AI-assisted reviews, not external
peer review.

\subsection{Complete coefficient bindings}

Two different complete four-port tuples are used.  They must not be
mixed.  The integer family uses the source-34 mixed-parity unrounded
tuple with $\beta=2^{1064921}-1$.  The positive family uses the compact
source-44 tuple with $\beta=2^{590069}-1$.  The tuple members are
$\beta,K_0,M_C,M_F$; $M_F$ already includes the source's prescribed
shift.  Reapplying that shift changes the equation.

The supplied compressed integer payloads have exact, documented byte
interpretations: after gzip decompression, bytes $b_0,\ldots,b_{m-1}$
represent $\sum_{j=0}^{m-1}b_j256^j$.  Thus the four ``ports'' in the
symbolic polynomial files are fixed ordinary integer coefficients,
not free parameters or uncomputed choices.  Their recipes, byte hashes,
and source pins are included.  No numeral is obtained by searching for
a contradiction or by deciding a consistency statement.

The four raw native payloads total 532,479 bytes for the integer family
and 295,388 bytes for the positive family.  This is storage for the
native inputs to the displayed polynomial recipes, not the size of a
fully expanded specialized polynomial.  The new affine restrictions
change small integer multipliers in the expanded equations.  Keeping
the native tuple unchanged therefore does not mean that every expanded
coefficient is unchanged.

\subsection{What the verification programs establish}

The root verification command authenticates the release manifest, makes
an isolated copy, and runs the component verification there.  It compares
regenerated literal equation artifacts with the delivered bytes.  Each
component uses the Python standard library, with checks that remain
active under \texttt{python3 -O}.  The checks cover:
\begin{enumerate}[leftmargin=1.5em]
\item exact affine substitutions from the pinned parent and deletion of
      precisely the intended coordinate and residual;
\item exact polynomial restoration maps for both integer companions;
\item counts, complete coordinate use, residual degrees, and nonzero
      monomial certificates of the claimed total degrees;
\item literal positive arithmetic DAGs, their live operation counts,
      cancellation-resolved degree identities, and fixed-numeral binding;
\item source-file identities, coefficient payload integrity, and the
      included finite native coefficient reconstruction.
\end{enumerate}

The integer audit also expands the sums of squares with the four fixed
ports temporarily treated as symbolic coefficients.  Its special
monomial certificates are independent of their large values.  The
positive degree audit gives explicit nonzero leading coefficients after
specialization to the actual positive $\beta$.  These lower bounds are
necessary: a degree propagated through a circuit is only an upper bound
until possible cancellation is addressed.

No verification command runs an unbounded refutation search, supplies a
zero of the consistency polynomial, or establishes unsolvability.
The polynomial coefficients are completely materialized, while the
existence equivalence is a theorem about all integer assignments.

\subsection{Using and integrating the archive}

The root directory contains this complete article as \texttt{report.tex}
and \texttt{report.pdf}, the component files under \texttt{integer/},
\texttt{positive/}, and \texttt{eleven/}, comparison provenance under \texttt{baseline/},
checks and receipts, licenses, and a root verification driver.  Paths
mentioned inside component proofs are relative to that component.
The cover and \texttt{README.md} identify the three main literal packets.
All residual systems in the integer JSON files define one equation by
taking their sum of squares.  The twelve- and thirteen-positive JSON
files include their sum-of-squares output nodes; the eleven-positive
file uses the proved unit-product finalizer.

Run \texttt{python3 verify.py} from the extracted root, or add
\texttt{--optimized} to exercise the checks with Python optimization
enabled.  Run \texttt{python3 build\_report.py} to compile the article
twice with \texttt{pdflatex}.  The mathematical tests need no network
access; the PDF build needs an ordinary TeX installation.  PDF metadata
can vary on recompilation, so the report rebuild is separate from exact
equation regeneration.

The directory is prepared as an additive contribution.  It can be placed
beside the earlier numbered packages at intake without changing their
delivered content or the active Lean equation.  The source snapshot is
fixed even if the repository subsequently advances.

\section{Further research questions}
\label{sec:future}

\paragraph{Can more sign witnesses be replaced by residues?}
The essential new theorem recovers a finite set of possible algebraic
values before using a congruence to select one.  A next target is to
classify which sign constraints in other low-degree Diophantine
compilers have this property.  A proposed deletion must prove the full
signed version of the earlier argument and preserve a canonical
completeness family.  Checking that the intended witnesses satisfy a
residue is insufficient to exclude unintended roots.

\paragraph{Can the common-sign guard be removed?}
The quartic still uses three coordinates to force
$\eta\zeta>0$.  They supply the sharp interval
$K(Z-1)<|c|<K(Z+1)$ needed for the rank comparison.  Replacing this
guard would save more variables than the single square witness removed
here, but the interval cannot simply be assumed after removing it.
A useful intermediate theorem would recover a suitable ratio interval
from retained congruences and norms alone.

\paragraph{Can a fourth root be transferred equally cheaply?}
The eleven-variable construction resolves the auxiliary integrality
obstacle through coprimality.  Eliminating the separately retained main
Pell carrier could improve another degree slice, but the three-root
rationality lemma does not extend by the same hypotheses alone.
For example,
\[
 10+\sqrt2+\sqrt8-\sqrt{18}=10
\]
has pairwise distinct radicands and a root larger than the sum of the
other three, yet three roots are irrational.  A four-root construction
needs control of proper signed subsums or another extraction mechanism,
as well as a new integrality and positive-completeness argument.  The
larger nested norms in the predecessors provide a concrete comparison.

\paragraph{Can the transport multiplier be made less expensive?}
In the twelve-variable equation, $W=1+(q-1)ikc$ pays for the positive
inverse gap.  The degree-$84$ leading form shows precisely where this
cost enters.  A smaller sufficient polynomial majorant for $L-f$,
compatible with the same congruence, could lower degree further.  Such
an estimate must be uniform in the cofinal strong-Pell witness, not
just true for a few small indices.

\paragraph{Which reductions are best after fixing a description language?}
Native tuple bytes, expanded monomials, circuit operations, and proof
complexity are different objectives.  A useful next comparison would
fix a concrete encoding of polynomial arithmetic DAGs and coefficients,
then compare total description length.  The present release supplies
enough literal data for that comparison; it does not infer it from
variable count alone.

\paragraph{Can a direct proof-language compiler become competitive?}
Relational NBG and other finitely presented equiconsistent systems
remain plausible alternatives to machine tableaux.  The missing
deliverable is a complete source predicate, a proved bridge to the
specific syntactic consistency statement, and a measured explicit
polynomial after all bounded-sequence and exponentiation devices are
compiled.  A smaller informal interpreter is not by itself a smaller
Diophantine equation.

\paragraph{What should be formalized next?}
The signed power-recovery lemma is a focused target: its inputs are
elementary Pell classifications, explicit magnitude bounds, and the
already used rank interfaces.  The two-root transfer lemma is an even
smaller independent target.  After those lemmas, the affine substitutions,
restoration maps, and literal DAG identities can be connected to the
repository's formal arithmetic syntax.  Completing the native semantic
and object-language internalization boundaries remains a separate task.

\begin{thebibliography}{9}
\bibitem{repo}
V. Reshetnikov and contributors.
\emph{ProveIt: DiophantineConsistency and delivered research packages}.
Snapshot \texttt{9cca523d9f27ad39d52993f4065cbd7e58f8bdd4}, 23 October 2026.
\url{https://github.com/VladimirReshetnikov/ProveIt/tree/9cca523d9f27ad39d52993f4065cbd7e58f8bdd4/Computability/DiophantineConsistency}.
Specific original paths and Git blob/SHA-256 identities are recorded in
the component \texttt{SOURCE\_PINS.json} files and root \texttt{BASELINE.json}.

\bibitem{active}
ProveIt contributors.
\emph{A Diophantine equation for syntactic consistency of ZFC}, active
research README at the same snapshot; retained as
\texttt{baseline/active\_research\_README.md}.
The active explicit equation and small-circuit formalization status are
distinct from the new equations in this release.

\bibitem{jones}
J. P. Jones.
Universal Diophantine equation.
\emph{Journal of Symbolic Logic} \textbf{47}(3), 549--571 (1982).
\url{https://doi.org/10.2307/2273588}.

\bibitem{carlmoroz}
M. Carl and B. Z. Moroz.
On a Diophantine representation of the predicate of provability.
\emph{Zapiski Nauchnykh Seminarov POMI} \textbf{407}, 77--104 (2012);
English translation, \emph{Journal of Mathematical Sciences}
\textbf{199}(1), 36--52 (2014).
\url{https://doi.org/10.1007/s10958-014-1830-2}.

\bibitem{fgo}
A. Formisano, I. Gavazzi, and E. G. Omodeo.
Diophantine Modeling of Provability in Algebraic Logic.
In \emph{CILC 2025}, CEUR Workshop Proceedings, vol. 4003, paper 13.
\url{https://ceur-ws.org/Vol-4003/paper13.pdf}.

\bibitem{sun}
Z.-W. Sun.
Ten unknowns for Hilbert's tenth problem over the integers.
arXiv:2609.01594v1 (2026).
\url{https://arxiv.org/abs/2609.01594v1}.
\end{thebibliography}
\end{document}
